% Book 3, Year 1, chapter 4 -- pointers and arrays.
% Provenance (never printed): sources/book3/BRIEFS.md, chapter 4.
% ledger: b1:ptr-size, b1:c-one-past-end, b1:qsort, b1:asan
\chapter{Pointers and Arrays}\label{ch:b1:pointers-and-arrays}

To read a whole number typed at the keyboard, a C program calls \texttt{scanf("\%d", \&n)}. Forget the ampersand, write
\texttt{scanf("\%d", n)}, and the \omterm{def:b1:compilation-pipeline:compiler}{compiler} warns while the program, with luck, crashes; without luck it writes the
number somewhere in memory and carries on. The \code{&} hands \code{scanf} the place where \code{n} lives, so that
it can store a value there: \cref{ch:b1:c-functions-and-stack} showed that a function receives copies, and a copy of
\code{n} would be of no use. This chapter is about such places, the addresses of memory, and about the values that
hold them, \omterm{def:b1:pointers-and-arrays:pointer}{pointers}. They are the most powerful feature of C and the source of its worst bugs, and the two are the
same thing.

\begin{recall}
Memory is a long row of bytes, each with a number, its \omterm{def:b1:pointers-and-arrays:pointer}{address}. A variable occupies a few consecutive bytes, from a
first \omterm{def:b1:pointers-and-arrays:pointer}{address} on: an \code{int} four bytes here, a \code{double} eight. The school volumes drew it as a row of
numbered boxes, and showed that a Python list holds references to its elements rather than the elements
themselves.
\end{recall}

\section{Addresses and pointers}\label{sec:b1:pointers-and-arrays:pointers}

\begin{definition}[Address, pointer, dereferencing]\label{def:b1:pointers-and-arrays:pointer}
The \emph{address}\index{address (of an object)} of an object is the address of its first byte; the operator
\code{&} gives it: \code{&n}. A \emph{pointer}\index{pointer} is a value that is an address together with the type
of what is stored there; a variable declared \code{int *p} holds a pointer to an \code{int}. The operator \code{*}
applied to a pointer, \code{*p}, designates the object at that address: reading \code{*p} reads it, assigning to
\code{*p} writes into it. Following a pointer in this way is called \emph{dereferencing}\index{dereferencing} it.
\end{definition}

\omcode{code/bachelor-1/04-pointers-and-arrays/listings/addr.c}{2}{16}{A variable and a \omterm{def:b1:pointers-and-arrays:pointer}{pointer} to
it}\label{lst:b1:pointers-and-arrays:addr}

\omoutput{code/bachelor-1/04-pointers-and-arrays/outputs/addr.txt}{Writing through \code{p} changed
\code{n}.}\label{out:b1:pointers-and-arrays:addr}

The type in \code{int *} matters twice: it tells the \omterm{def:b1:compilation-pipeline:compiler}{compiler} how many bytes \code{*p} covers and how to read them,
and it lets the \omterm{def:b1:compilation-pipeline:compiler}{compiler} refuse \code{double *q = &n;}. A \omterm{def:b1:pointers-and-arrays:pointer}{pointer} itself occupies 8 bytes on this machine, whatever
it points to, because every \omterm{def:b1:pointers-and-arrays:pointer}{address} of a 64-bit process fits in 64 bits. The \code{*} has two roles in the text: in
a declaration it says ``\omterm{def:b1:pointers-and-arrays:pointer}{pointer} to''; in an expression it dereferences.

\begin{omfigure}
\begin{tikzpicture}[c/.style={draw, minimum width=16mm, minimum height=7mm, font=\small\ttfamily},
                    >={Stealth[length=2mm]}]
\foreach \k/\a in {0/1000, 1/1004, 2/1008, 3/1012} {
  \node[font=\scriptsize\ttfamily, omCode] at (\k*1.6,0.65) {\a};
}
\node[c, fill=omDef!8] (n) at (0,0) {7};
\node[c] at (1.6,0) {};
\node[c, fill=omThm!8, minimum width=32mm] (p) at (4.0,0) {1000};
\node[anchor=north west, font=\small] at (n.south west) {\code{n}};
\node[anchor=north east, font=\small] at (p.south east) {\code{p}};
\draw[->, thick, omThm] (p.south) to[out=-110, in=-70] node[below, font=\scriptsize\rmfamily] {\code{*p} is
  \code{n}} (n.south);
\end{tikzpicture}
\par\smallskip{\small Memory after \code{*p = 7}, with made-up addresses (the real ones change from run to run).
\code{n} occupies the four bytes from \omterm{def:b1:pointers-and-arrays:pointer}{address} 1000; \code{p}, eight bytes from 1008, holds the number 1000. Writing
\code{*p} writes into the bytes at that number, which are \code{n}'s.}
\end{omfigure}

\begin{definition}[Null pointer]\label{def:b1:pointers-and-arrays:null}
The \emph{null pointer}\index{null pointer} of a \omterm{def:b1:pointers-and-arrays:pointer}{pointer} type, written \code{NULL} (or 0), compares unequal to the
\omterm{def:b1:pointers-and-arrays:pointer}{address} of every object and function. It means ``points to nothing'': a function that finds nothing returns it, and
a \omterm{def:b1:pointers-and-arrays:pointer}{pointer} variable not yet given a target should hold it. \omterm{def:b1:pointers-and-arrays:pointer}{Dereferencing} the null \omterm{def:b1:pointers-and-arrays:pointer}{pointer} is \omterm{def:b1:c-types-and-control:ub}{undefined behaviour}.
\end{definition}

On this system the first page of memory is never mapped, and \omterm{def:b1:pointers-and-arrays:pointer}{dereferencing} \code{NULL} stops the program at once, the
best outcome an \omterm{def:b1:c-types-and-control:ub}{undefined behaviour} can have (\cref{sec:b1:pointers-and-arrays:oob}).

\section{A pointer passed to change the caller's data}\label{sec:b1:pointers-and-arrays:passing}

\omterm{def:b1:c-functions-and-stack:by-value}{Passing by value} still holds: a function receives copies. But the copy of an \omterm{def:b1:pointers-and-arrays:pointer}{address} is the same \omterm{def:b1:pointers-and-arrays:pointer}{address}, and through
it the function reaches the caller's variable. The failed \code{swap} of \cref{ch:b1:c-functions-and-stack} works
once it receives addresses.

\omcode{code/bachelor-1/04-pointers-and-arrays/listings/swap_ptr.c}{2}{17}{A swap that
swaps}\label{lst:b1:pointers-and-arrays:swap}

\omoutput{code/bachelor-1/04-pointers-and-arrays/outputs/swap_ptr.txt}{This time \code{x} and \code{y} are
exchanged.}

\begin{omfigure}
\begin{tikzpicture}[f/.style={draw, minimum width=36mm, minimum height=6mm, font=\small, anchor=north},
                    >={Stealth[length=2mm]}]
\node[f, fill=omCode!6] (x) at (0,0) {\code{x = 2}};
\node[f, fill=omCode!6] (y) at (0,-0.6) {\code{y = 1}};
\node[f, fill=omThm!10] (r) at (0,-1.2) {return address};
\node[f, fill=omDef!8] (b) at (0,-1.8) {\code{b = &y}};
\node[f, fill=omDef!8] (a) at (0,-2.4) {\code{a = &x}};
\node[f, fill=omDef!8] (t) at (0,-3.0) {\code{t = 1}};
\draw[->, thick, omDef] (a.east) to[out=0, in=0, looseness=1.4] (x.east);
\draw[->, thick, omDef] (b.east) to[out=0, in=0] (y.east);
\draw[decorate, decoration={brace, amplitude=4pt, mirror}] (-1.95,0) -- (-1.95,-1.2)
  node[midway, left=5pt, font=\small] {frame of \code{main}};
\draw[decorate, decoration={brace, amplitude=4pt, mirror}] (-1.95,-1.2) -- (-1.95,-3.6)
  node[midway, left=5pt, font=\small] {frame of \code{swap}};
\end{tikzpicture}
\par\smallskip{\small The frames at the end of the fixed \code{swap}. Its parameters hold the addresses of
\code{main}'s variables: writing \code{*a} and \code{*b} writes into \code{main}'s frame. This is the only way a
called function reaches the frame of its caller.}
\end{omfigure}

The same device gives a function several results: it returns one, or none, and writes the others through \omterm{def:b1:pointers-and-arrays:pointer}{pointer}
parameters supplied by the caller. \code{scanf} works this way: it returns the number of values it converted, and
stores each value at the \omterm{def:b1:pointers-and-arrays:pointer}{address} it was given.

\omcode{code/bachelor-1/04-pointers-and-arrays/listings/divmod.c}{2}{16}{Two results through two
\omterm{def:b1:pointers-and-arrays:pointer}{pointers}}\label{lst:b1:pointers-and-arrays:divmod}

\omoutput{code/bachelor-1/04-pointers-and-arrays/outputs/divmod.txt}{$2026 = 7 \times 289 + 3$.}

\omcode{code/bachelor-1/04-pointers-and-arrays/listings/read_number.c}{2}{13}{Reading a number: \code{scanf} stores
it at \code{&n}}\label{lst:b1:pointers-and-arrays:read}

\omoutput{code/bachelor-1/04-pointers-and-arrays/outputs/read_number.txt}{With 45 on standard input.}

A \omterm{def:b1:pointers-and-arrays:pointer}{pointer} parameter that the function only reads from is declared \code{const int *p}: the function receives a
read-only view and the \omterm{def:b1:compilation-pipeline:compiler}{compiler} refuses \code{*p = 0}. The \code{const} documents the contract and is checked: a
reader of the \omterm{def:b1:c-functions-and-stack:declaration}{prototype} \code{int min_max(const int *a, size_t n, int *min, int *max)} knows that the \omterm{def:b1:pointers-and-arrays:array}{array} is
read and that \code{*min} and \code{*max} are written.

\section{Arrays and pointer arithmetic}\label{sec:b1:pointers-and-arrays:arrays}

\begin{definition}[Array, array-to-pointer conversion]\label{def:b1:pointers-and-arrays:array}
An \emph{array}\index{array (C)} of C, declared \code{int a[5]}, is a block of a fixed number of elements of one type,
stored contiguously in memory; its size is fixed when it is created. In almost every expression, the name of an array
is converted into a \omterm{def:b1:pointers-and-arrays:pointer}{pointer} to its first element: this is the \emph{array-to-pointer
conversion}\index{array-to-pointer conversion}. The exceptions are \code{sizeof a} and \code{&a}, where \code{a}
stays the whole array.
\end{definition}

\begin{definition}[Pointer arithmetic]\label{def:b1:pointers-and-arrays:arithmetic}
If \code{p} points to element $i$ of an \omterm{def:b1:pointers-and-arrays:array}{array}, \code{p + k} points to element $i + k$: \emph{pointer
arithmetic}\index{pointer arithmetic} counts in elements, not bytes, so the \omterm{def:b1:pointers-and-arrays:pointer}{address} grows by $k$ times the size of an
element. The result must point into the \omterm{def:b1:pointers-and-arrays:array}{array} or exactly one element past its end. The difference \code{q - p} of two
\omterm{def:b1:pointers-and-arrays:pointer}{pointers} into the same \omterm{def:b1:pointers-and-arrays:array}{array} is the number of elements between them, of type \code{ptrdiff_t}. By definition
\code{a[i]} means \code{*(a + i)}.
\end{definition}

\omcode{code/bachelor-1/04-pointers-and-arrays/listings/array_layout.c}{2}{14}{The layout of an
\omterm{def:b1:pointers-and-arrays:array}{array}}\label{lst:b1:pointers-and-arrays:layout}

\omoutput{code/bachelor-1/04-pointers-and-arrays/outputs/array_layout.txt}{Elements four bytes apart; \code{a + 5}
is one past the end.}\label{out:b1:pointers-and-arrays:layout}

The \omterm{def:b1:pointers-and-arrays:pointer}{pointer} \code{a + 5} is allowed to exist, to be compared and subtracted, but not to be dereferenced: it is the
end mark of every loop of the form \code{for (p = a; p < a + n; p++)}. A \omterm{def:b1:pointers-and-arrays:pointer}{pointer} further away, or before the start,
is \omterm{def:b1:c-types-and-control:ub}{undefined behaviour} even if it is never dereferenced.

\begin{omfigure}
\begin{tikzpicture}[c/.style={draw, minimum width=15mm, minimum height=7mm, font=\small\ttfamily},
                    >={Stealth[length=2mm]}]
\foreach \k/\v in {0/10, 1/20, 2/30, 3/40, 4/50} {
  \node[c, fill=omDef!8] (e\k) at (\k*1.5,0) {\v};
  \node[font=\scriptsize\ttfamily, omCode] at (\k*1.5,0.62) {+\the\numexpr\k*4\relax};
  \node[font=\small\ttfamily] (l\k) at (\k*1.5,-1.2) {a+\k};
  \draw[->] (l\k) -- (e\k);
}
\node[c, dashed, draw=omThm] (e5) at (7.5,0) {};
\node[font=\scriptsize\ttfamily, omCode] at (7.5,0.62) {+20};
\node[font=\small\ttfamily, omThm] (l5) at (7.5,-1.2) {a+5};
\draw[->, omThm] (l5) -- (e5);
\node[font=\scriptsize, omThm, align=center] at (7.5,-1.85) {one past the end:\\compare, never read};
\node[font=\scriptsize, align=left, anchor=west] at (-0.8,1.2) {byte offset from \code{a}};
\end{tikzpicture}
\par\smallskip{\small An \omterm{def:b1:pointers-and-arrays:array}{array} of five \code{int}. The \omterm{def:b1:pointers-and-arrays:pointer}{pointer} \code{a + k} designates element $k$, $4k$ bytes from the
start; \code{a[k]} and \code{*(a + k)} are the same element. The position after the last element can be pointed to
but holds nothing of the \omterm{def:b1:pointers-and-arrays:array}{array}.}
\end{omfigure}

Since an \omterm{def:b1:pointers-and-arrays:array}{array} name becomes a \omterm{def:b1:pointers-and-arrays:pointer}{pointer}, an \omterm{def:b1:pointers-and-arrays:array}{array} is never passed to a function: its \omterm{def:b1:pointers-and-arrays:pointer}{address} is. A parameter written
\code{int a[]} is a \omterm{def:b1:pointers-and-arrays:pointer}{pointer} \code{int *a}, the function cannot know how many elements follow, and the length must
travel in a separate parameter. \code{sizeof} on such a parameter gives the size of a \omterm{def:b1:pointers-and-arrays:pointer}{pointer}; GCC warns.

\omcode{code/bachelor-1/04-pointers-and-arrays/listings/decay.c}{2}{19}{An \omterm{def:b1:pointers-and-arrays:array}{array} parameter is a
\omterm{def:b1:pointers-and-arrays:pointer}{pointer}}\label{lst:b1:pointers-and-arrays:decay}

\omoutput{code/bachelor-1/04-pointers-and-arrays/outputs/decay.txt}{48 bytes in \code{main}, 8 in
\code{sum}.}\label{out:b1:pointers-and-arrays:decay}

This is also why a function can change the elements of an \omterm{def:b1:pointers-and-arrays:array}{array} it receives, with no \code{&} at the call: the
\omterm{def:b1:pointers-and-arrays:array}{array}'s \omterm{def:b1:pointers-and-arrays:pointer}{address} is passed \omterm{def:b1:c-functions-and-stack:by-value}{by value}, and the elements are reached through it.

\begin{definition}[Row-major order]\label{def:b1:pointers-and-arrays:row-major}
A two-dimensional \omterm{def:b1:pointers-and-arrays:array}{array} \code{int g[h][w]} is an \omterm{def:b1:pointers-and-arrays:array}{array} of $h$ rows, each an \omterm{def:b1:pointers-and-arrays:array}{array} of $w$ elements, stored one row
after the other: this is \emph{row-major order}\index{row-major order}. Element \code{g[i][j]} is element
$i \cdot w + j$ of the block of $h \cdot w$ elements.
\end{definition}

\omcode{code/bachelor-1/04-pointers-and-arrays/listings/rowmajor.c}{2}{16}{Twelve \code{int} in a
row}\label{lst:b1:pointers-and-arrays:rowmajor}

\omoutput{code/bachelor-1/04-pointers-and-arrays/outputs/rowmajor.txt}{\code{g[2][1]} is element $2 \cdot 4 + 1 = 9$.}

Every \omterm{def:b1:compilation-pipeline:compiler}{compiler} lays the rows out this way and lets \code{flat} walk the twelve elements. Strictly, though, the
standard lets a \omterm{def:b1:pointers-and-arrays:pointer}{pointer} derived from the row \code{g[0]} range only over that row and one past its end; walking
further, or subtracting a \omterm{def:b1:pointers-and-arrays:pointer}{pointer} into row 2, is formally undefined. Portable code indexes \code{g[i][j]}.

\begin{omfigure}
\begin{tikzpicture}[c/.style={draw, minimum width=8.5mm, minimum height=6.5mm, font=\scriptsize\ttfamily}]
\foreach \i in {0,1,2} \foreach \j in {0,1,2,3} {
  \pgfmathtruncatemacro{\v}{10*\i+\j}
  \node[c, fill=omDef!\the\numexpr 6+6*\i\relax] at (\j*0.85, -\i*0.65) {\v};
}
\node[font=\small, anchor=east] at (-0.6,-0.65) {\code{g[3][4]}};
\foreach \k in {0,...,11} {
  \pgfmathtruncatemacro{\i}{\k/4}
  \pgfmathtruncatemacro{\j}{mod(\k,4)}
  \pgfmathtruncatemacro{\v}{10*\i+\j}
  \node[c, fill=omDef!\the\numexpr 6+6*\i\relax] at (\k*0.85+0.0, -2.6) {\v};
  \node[font=\tiny\ttfamily, omCode] at (\k*0.85, -3.15) {\k};
}
\node[font=\small, anchor=east] at (-0.6,-2.6) {memory};
\draw[->, >={Stealth[length=2mm]}] (1.3,-1.6) -- node[right, font=\scriptsize] {row after row} (1.3,-2.2);
\node[font=\scriptsize, align=left, anchor=west] at (3.9,-0.65) {element $(i, j)$ is at\\index $i \cdot 4 + j$};
\end{tikzpicture}
\par\smallskip{\small The \omterm{def:b1:pointers-and-arrays:array}{array} \code{g[3][4]} of \cref{lst:b1:pointers-and-arrays:rowmajor}, as written in the source
and as stored. The shade marks the row. Walking a row is a stride of 1 in memory; walking a column, a stride of 4.}
\end{omfigure}

\section{Out of bounds}\label{sec:b1:pointers-and-arrays:oob}

\begin{definition}[Out-of-bounds access]\label{def:b1:pointers-and-arrays:oob}
An \emph{out-of-bounds access}\index{out-of-bounds access} reads or writes through a \omterm{def:b1:pointers-and-arrays:pointer}{pointer}, or an index, that does
not designate an element of the \omterm{def:b1:pointers-and-arrays:array}{array} it was derived from. C does not check indices: an out-of-bounds access is
\omterm{def:b1:c-types-and-control:ub}{undefined behaviour}, and in practice it reads or writes whatever lies at the computed \omterm{def:b1:pointers-and-arrays:pointer}{address}.
\end{definition}

The commonest one is a loop that runs once too many. In the program below, the fifth write of the loop meant for
\code{marks} lands in \code{bonus}, which on this build is stored just after it; no message, no crash, and a wrong
result far from the cause.

\omcode{code/bachelor-1/04-pointers-and-arrays/listings/neighbour.c}{3}{13}{One write too
many}\label{lst:b1:pointers-and-arrays:neighbour}

\omoutput{code/bachelor-1/04-pointers-and-arrays/outputs/neighbour.txt}{\code{bonus[0]} was changed by a loop that
never names it.}\label{out:b1:pointers-and-arrays:neighbour}

The \omterm{def:b1:pointers-and-arrays:pointer}{address} sanitizer, enabled by \code{-fsanitize=address}, places a forbidden zone around every \omterm{def:b1:pointers-and-arrays:array}{array} and checks
every access; \code{-g} lets it report source lines. The same program, built that way, stops at the bad write and
explains it.

\omoutput[1-6,14-16]{code/bachelor-1/04-pointers-and-arrays/outputs/neighbour_asan.txt}{The first lines of the sanitizer's
report (the rest is a map of the forbidden zones).}\label{out:b1:pointers-and-arrays:asan}

Read the report from the top: the kind of error (a buffer on the stack overflowed), the access (a write of 4 bytes),
the place in the source (line 10, \code{marks[i] = 0}), then the variable it overflowed: \code{marks} occupies
offsets 32 to 48 of \code{main}'s frame and the write hit offset 48. As with the undefined-behaviour sanitizer, it
can only report the accesses that the run performs.

A \omterm{def:b1:pointers-and-arrays:null}{null pointer} dereference needs no sanitizer: the system stops the program.

\omcode{code/bachelor-1/04-pointers-and-arrays/listings/nullderef.c}{2}{13}{Reading through the \omterm{def:b1:pointers-and-arrays:null}{null
pointer}}\label{lst:b1:pointers-and-arrays:null}

\omoutput{code/bachelor-1/04-pointers-and-arrays/outputs/nullderef.txt}{Status 139, as for the \omterm{def:b1:c-functions-and-stack:stack-overflow}{stack overflow} of
\cref{ch:b1:c-functions-and-stack}: a segmentation fault.}\label{out:b1:pointers-and-arrays:null}

The status 139 reports a \emph{segmentation fault}: an access to an \omterm{def:b1:pointers-and-arrays:pointer}{address} that the process has no right to. It is
the system's protection, not the language's; an \omterm{def:b1:pointers-and-arrays:oob}{out-of-bounds access} that stays inside memory the process owns, as
in \code{neighbour.c}, is not caught at all.

\begin{safety}
Every \omterm{def:b1:pointers-and-arrays:array}{array} access in C is a promise by the programmer that the index is in range. Keep the length beside the
\omterm{def:b1:pointers-and-arrays:array}{array}, write loops as \code{i < n}, never \code{i <= n}, and run the tests with \code{-fsanitize=address}. Most of
the security flaws in C programs of the last decades started as an out-of-bounds write.
\end{safety}

\section{Pointers to functions}\label{sec:b1:pointers-and-arrays:functions}

\begin{definition}[Function pointer]\label{def:b1:pointers-and-arrays:function-pointer}
A \emph{function pointer}\index{function pointer} holds the \omterm{def:b1:pointers-and-arrays:pointer}{address} of a function together with its type: for
example \code{int (*f)(int, int)} declares \code{f} as a \omterm{def:b1:pointers-and-arrays:pointer}{pointer} to a function taking two \code{int} and returning
an \code{int}. The name of a function converts to a \omterm{def:b1:pointers-and-arrays:pointer}{pointer} to it, and a call \code{f(6, 7)} through the \omterm{def:b1:pointers-and-arrays:pointer}{pointer}
calls the function it points to.
\end{definition}

\omcode{code/bachelor-1/04-pointers-and-arrays/listings/ops_table.c}{2}{17}{A table of
operations}\label{lst:b1:pointers-and-arrays:ops}

\omoutput{code/bachelor-1/04-pointers-and-arrays/outputs/ops_table.txt}{The operation is chosen by an index at run
time.}

A \omterm{def:b1:pointers-and-arrays:function-pointer}{function pointer} lets a function receive part of its algorithm as an argument. The C library's sorting function,
\code{qsort}, sorts an \omterm{def:b1:pointers-and-arrays:array}{array} of any type: it receives the \omterm{def:b1:pointers-and-arrays:array}{array}'s \omterm{def:b1:pointers-and-arrays:pointer}{address}, the number of elements, the size of one
element, and a comparison function that it calls with the addresses of two elements. Since \code{qsort} cannot know
the type, those addresses arrive as \code{const void *}.

\begin{definition}[Generic pointer]\label{def:b1:pointers-and-arrays:generic}
A \omterm{def:b1:pointers-and-arrays:pointer}{pointer} to \code{void}, or \emph{generic pointer}\index{generic pointer}, holds an \omterm{def:b1:pointers-and-arrays:pointer}{address} without saying what is
stored there. Any object \omterm{def:b1:pointers-and-arrays:pointer}{pointer} converts to it and back without loss, with no \omterm{def:b1:c-types-and-control:cast}{cast} in C. It cannot be dereferenced
nor used in \omterm{def:b1:pointers-and-arrays:arithmetic}{pointer arithmetic}: it must first be converted to a \omterm{def:b1:pointers-and-arrays:pointer}{pointer} to a real type.
\end{definition}

\omcode{code/bachelor-1/04-pointers-and-arrays/listings/sort_people.c}{6}{36}{Three orders for \code{qsort}, three
comparators}\label{lst:b1:pointers-and-arrays:qsort}

\omoutput{code/bachelor-1/04-pointers-and-arrays/outputs/sort_people.txt}{Birth years of eight people of the town,
then their first names, in three orders.}\label{out:b1:pointers-and-arrays:qsort}

A comparator returns a negative number, zero or a positive number as its first argument comes before, with or after
the second. The expression \code{(a > b) - (a < b)} computes exactly that without the subtraction \code{a - b},
which could overflow. To sort the names, the elements are \omterm{def:b1:pointers-and-arrays:pointer}{pointers} to characters; the comparator receives the
addresses of two of them, hence the two levels of \omterm{def:b1:pointers-and-arrays:pointer}{pointer} in \code{*(const char *const *)x}. Strings are the
subject of \cref{ch:b1:c-strings-structs-files}, and \cref{ch:b1:dynamic-memory} uses \code{void *} at every step.

\begin{history}[Pointers in the hardware]
C's \omterm{def:b1:pointers-and-arrays:pointer}{pointers} are a thin layer over the machine. The PDP-11, on which C grew up, addressed memory byte by byte and had
instructions that read through a register and then advanced it by the size of the item read, which made
\code{*p++} a single instruction. BCPL and B, C's ancestors, had a single type, the machine word; C added types to
\omterm{def:b1:pointers-and-arrays:pointer}{pointers} so that \code{p + 1} could step by the size of the element, whatever it was.
\end{history}

\section{Exercises}\label{sec:b1:pointers-and-arrays:exercises}

\begin{exercise}[$\star$]\label{exo:b1:pointers-and-arrays:1}
Predict the line printed, following the \omterm{def:b1:pointers-and-arrays:pointer}{pointers} one statement at a time.
\omcode{code/bachelor-1/04-pointers-and-arrays/listings/exo_ptrs.c}{2}{15}{Follow the \omterm{def:b1:pointers-and-arrays:pointer}{pointers}}
\end{exercise}

\begin{exercise}[$\star$]\label{exo:b1:pointers-and-arrays:2}
Draw the memory at the end of \code{main}, as in the figures of this chapter: the \omterm{def:b1:pointers-and-arrays:array}{array}, the three \omterm{def:b1:pointers-and-arrays:pointer}{pointers}, and an
arrow for each \omterm{def:b1:pointers-and-arrays:pointer}{pointer}. What does \code{main} return?
\omcode{code/bachelor-1/04-pointers-and-arrays/listings/exo_diagram.c}{2}{11}{A \omterm{def:b1:pointers-and-arrays:pointer}{pointer} to a \omterm{def:b1:pointers-and-arrays:pointer}{pointer}}
\end{exercise}

\begin{exercise}[$\star$]\label{exo:b1:pointers-and-arrays:3}
Predict the four values printed (a \code{double} occupies 8 bytes).
\omcode{code/bachelor-1/04-pointers-and-arrays/listings/exo_diff.c}{2}{10}{\omterm{def:b1:pointers-and-arrays:pointer}{Pointers} into an \omterm{def:b1:pointers-and-arrays:array}{array} of \code{double}}
\end{exercise}

\begin{exercise}[$\star$]\label{exo:b1:pointers-and-arrays:4}
What does this program print, and why does \code{length} not compute the length? Rewrite \code{length}'s caller so
that the length is computed where it can be.
\omcode{code/bachelor-1/04-pointers-and-arrays/listings/exo_sizeof.c}{2}{14}{The length, computed twice}
\end{exercise}

\begin{exercise}[$\star\star$]\label{exo:b1:pointers-and-arrays:5}
Explain the warning below. What would \code{scanf} do with the value it received if the program ran, and why might it
crash or might not?
\omcode{code/bachelor-1/04-pointers-and-arrays/listings/exo_scanf.c}{2}{10}{Reading a number, the wrong way}
\omoutput{code/bachelor-1/04-pointers-and-arrays/outputs/exo_scanf.txt}{The \omterm{def:b1:compilation-pipeline:compiler}{compiler}'s warning.}
\end{exercise}

\begin{exercise}[$\star\star$]\label{exo:b1:pointers-and-arrays:6}
The program below was built with the \omterm{def:b1:pointers-and-arrays:pointer}{address} sanitizer. Read the report, give the bug and its fix, and say what the
program would have printed without the sanitizer, in the best case.
\omcode{code/bachelor-1/04-pointers-and-arrays/listings/exo_loop.c}{3}{15}{The average rainfall}
\omoutput[1-6,11-12]{code/bachelor-1/04-pointers-and-arrays/outputs/exo_loop_asan.txt}{The first lines of the report.}
\end{exercise}

\begin{exercise}[$\star\star$]\label{exo:b1:pointers-and-arrays:7}
Write \code{int min_max(const int *a, size_t n, int *min, int *max)}, which stores the smallest and the largest
element of \code{a[0..n-1]} and returns 0, or returns $-1$ without storing anything when \code{n} is 0.
\end{exercise}

\begin{exercise}[$\star\star$]\label{exo:b1:pointers-and-arrays:8}
Write \code{size_t reverse(int *a, size_t n)}, which reverses an \omterm{def:b1:pointers-and-arrays:array}{array} in place with two \omterm{def:b1:pointers-and-arrays:pointer}{pointers}, one moving
forwards from the start and one backwards from the end, and returns the number of swaps it made.
\end{exercise}

\begin{exercise}[$\star\star$]\label{exo:b1:pointers-and-arrays:9}
Predict what the program prints, and explain from the definition of \code{a[i]} why \code{2[a]} compiles.
\omcode{code/bachelor-1/04-pointers-and-arrays/listings/exo_index.c}{2}{10}{\code{a[2]} and \code{2[a]}}
\end{exercise}

\begin{exercise}[$\star\star\star$]\label{exo:b1:pointers-and-arrays:10}
Prove that the \omterm{def:b1:pointers-and-arrays:pointer}{two-pointer} reversal of \cref{exo:b1:pointers-and-arrays:8} is correct and makes exactly
$\lfloor n/2 \rfloor$ swaps. State the invariant of its loop in terms of \code{left - a} and \code{right - a}.
\end{exercise}

\begin{exercise}[$\star\star\star$]\label{exo:b1:pointers-and-arrays:11}
Write \code{generic_max(const void *base, size_t n, size_t size, int (*cmp)(const void *, const void *))}, which
returns the \omterm{def:b1:pointers-and-arrays:pointer}{address} of a largest element of an \omterm{def:b1:pointers-and-arrays:array}{array} of any type, compared by \code{cmp} with the same contract as
for \code{qsort}, or \code{NULL} when \code{n} is 0. Why must it convert \code{base} to a \omterm{def:b1:pointers-and-arrays:pointer}{pointer} to \code{char}?
\end{exercise}

\begin{exercise}[$\star\star\star$]\label{exo:b1:pointers-and-arrays:12}
In an \omterm{def:b1:pointers-and-arrays:array}{array} of $n$ elements, the valid \omterm{def:b1:pointers-and-arrays:pointer}{pointers} are \code{a + 0} to \code{a + n} (one past the end included). Count
the pairs $(p, q)$ of valid \omterm{def:b1:pointers-and-arrays:pointer}{pointers} with $p \leq q$, and prove your count. What does each such pair describe?
\end{exercise}

\section{Problem: A spreadsheet in a flat array}\label{sec:b1:pointers-and-arrays:problem}

\begin{problem}[{Weekend problem --- A spreadsheet in a flat array}]\label{pb:b1:pointers-and-arrays:1}
A small spreadsheet program stores a grid of $h$ rows and $w$ columns of whole numbers in one block of $h \cdot w$
\code{int}, row after row. The module \code{grid.c} computes the \omterm{def:b1:pointers-and-arrays:pointer}{address} of a cell, sums along a row and a column, and
transposes a square grid in place; \code{sheet.c} tries it on a $4 \times 4$ grid. A colleague's first transpose,
\code{transpose_bug.c}, was built with the \omterm{def:b1:pointers-and-arrays:pointer}{address} sanitizer.
\omcode{code/bachelor-1/04-pointers-and-arrays/listings/grid.c}{5}{37}{The module: the index map, two walks, the
transpose}\label{lst:b1:pointers-and-arrays:grid}
\omoutput{code/bachelor-1/04-pointers-and-arrays/outputs/sheet.txt}{The grid, its sums, its transpose, and the swaps
for $n = 1$ to 8.}\label{out:b1:pointers-and-arrays:sheet}
\omcode{code/bachelor-1/04-pointers-and-arrays/listings/transpose_bug.c}{5}{13}{The colleague's
transpose}\label{lst:b1:pointers-and-arrays:bug}
\omoutput[1-7,15-16]{code/bachelor-1/04-pointers-and-arrays/outputs/transpose_asan.txt}{The first lines of the
sanitizer's report.}\label{out:b1:pointers-and-arrays:bug}

\needspace{4\baselineskip}
\medskip\noindent\textbf{Part I --- The index map.}
\begin{enumerate}
\item Let $f(i, j) = i w + j$ for $0 \leq i < h$, $0 \leq j < w$. Compute $f(2, 1)$ and $f(3, 3)$ for $w = 4$.
\item Prove that $f$ takes its values in $[0, hw)$.
\item Prove that $f$ is injective. (Hint: if $iw + j = i'w + j'$, look at the remainder of the division by $w$.)
\item Deduce that $f$ is a bijection onto $[0, hw)$, and give its inverse with \code{/} and \texttt{\%}.
\item Why does \code{cell} return \code{g + i * w + j} and not \code{g + (i * w + j) * sizeof(int)}?
\item A grid of $h = 50\,000$ rows and $w = 50\,000$ columns: can \code{i * w + j} be computed in \code{int}? What
type should the index have?
\end{enumerate}

\needspace{4\baselineskip}
\medskip\noindent\textbf{Part II --- Walking with pointers.}
\begin{enumerate}[resume]
\item In \code{row_sum}, which cells does \code{p} visit, and what is the last value of \code{p}, which is never
dereferenced?
\item In \code{col_sum}, by how many elements does \code{p} advance at each step? By how many bytes?
\item Prove that \code{col_sum} visits exactly the $h$ cells of column $j$.
\item The last value of \code{p} in \code{col_sum} is \code{g + h * w + j}. Is it allowed to compute it? Rewrite the
loop with an index so that no \omterm{def:b1:pointers-and-arrays:pointer}{pointer} beyond the end is formed.
\item Check the printed sums: row 2 and column 3 of the grid \code{10 * i + j}.
\item Give, for the grid \code{10 * i + j} with $h = w = n$, the formulas of row sum $i$ and column sum $j$.
\end{enumerate}

\needspace{4\baselineskip}
\medskip\noindent\textbf{Part III --- The transpose in place.}
\begin{enumerate}[resume]
\item Show that exchanging \code{g[i * n + j]} and \code{g[j * n + i]} for every pair with $i < j$ transposes the grid.
\item Why must the loop not also run over the pairs with $i > j$? What would the grid be after it?
\item Count the swaps: prove that there are exactly $n(n-1)/2$.
\item Check that count against the output for $n = 1$ to 8.
\item Show that the diagonal cells $(i, i)$ never move, and that each other cell moves exactly once.
\item How many swaps would a transpose into a second grid need, and how much extra memory? Compare.
\end{enumerate}

\needspace{4\baselineskip}
\medskip\noindent\textbf{Part IV --- The colleague's bug.}
\begin{enumerate}[resume]
\item For $n = 3$ and $i = 0$, list the pairs of indices that the buggy loop exchanges.
\item Which of them is out of the \omterm{def:b1:pointers-and-arrays:array}{array}, and why does the sanitizer report a read of 4 bytes at offset 84 of the
frame, given that \code{g} occupies offsets 48 to 84?
\item Without the sanitizer, which two cells does the exchange with $j = n$ touch, for $i < n - 1$? Which of them lies inside the \omterm{def:b1:pointers-and-arrays:array}{array}, and what does it hold?
\item For $n = 1$, which exchange does the buggy loop make? Explain why a test on a $1 \times 1$ grid, without the sanitizer, would not notice the bug.
\item Which line of the report lets you find the faulty line of the source, and which tells you that the caller is
\code{main}?
\item State the named result of the problem: \emph{the in-place transpose makes exactly $n(n-1)/2$ swaps}, with the
range of pairs that achieves it and the counts measured.
\end{enumerate}
\end{problem}
